% Shared TikZ styles for the Odin III paper figures (IMAGE/fig_*.tex).
% Needs: \usepackage{tikz} and
%   \usetikzlibrary{positioning,fit,arrows.meta,shapes.geometric,backgrounds,calc,matrix}
% loaded first (odin3.tex and IMAGE/figures.tex both do). Input this file once, in the preamble.
%
% Palette: the greens and browns of a tree and its greenery, with black for Odin's raven as the
% strong ink. Green = ColorBrewer "Greens", brown = a warm earth ramp. Text is dark on the light
% fills and white only on the dark browns and the black. Build status is coded twice, by fill and
% by border, so it survives a black-and-white print:
%   built        mid-green fill,  solid dark-green border
%   in progress  pale-green fill, solid green border
%   specified    white fill,      dashed green border
%   planned      white fill,      dotted brown border

% ---- colours ------------------------------------------------------------------------------
\definecolor{o3ink}{HTML}{111111}      % text and arrows: the raven's black
\definecolor{o3sub}{HTML}{4E342E}      % secondary text (sub-project tags, notes): dark brown
\definecolor{o3grn9}{HTML}{00441B}
\definecolor{o3grn8}{HTML}{006D2C}
\definecolor{o3grn7}{HTML}{238B45}
\definecolor{o3grn5}{HTML}{74C476}
\definecolor{o3grn4}{HTML}{A1D99B}
\definecolor{o3grn3}{HTML}{C7E9C0}
\definecolor{o3grn2}{HTML}{E5F5E0}
\definecolor{o3grn1}{HTML}{F7FCF5}
\definecolor{o3brn9}{HTML}{3E2723}     % root system, dark brown (white text only)
\definecolor{o3brn8}{HTML}{5D4037}
\definecolor{o3brn7}{HTML}{4E342E}
\definecolor{o3brn6}{HTML}{6D4C41}     % trunk
\definecolor{o3brn5}{HTML}{8D6E63}
\definecolor{o3brn4}{HTML}{C9B6A4}     % branches, lighter brown
\definecolor{o3brn3}{HTML}{E0D3C6}
\definecolor{o3brn2}{HTML}{F1E9DF}
\definecolor{o3brn1}{HTML}{F8F4EE}

% ---- text helpers -------------------------------------------------------------------------
% A second, smaller line in a node: \odinsub{2C}.
\newcommand{\odinsub}[1]{{\scriptsize\color{o3sub}#1}}
% Code names inside figures (the paper's \code is not defined in the standalone build).
\newcommand{\odintt}[1]{\texttt{#1}}

% Scale a figure down only if it is wider than #1 (it never scales up). Each figure is drawn
% at its print size, so this is a safety net; a scaled figure writes "ODINFIT" to the log.
\newsavebox{\odinfigbox}
\newcommand{\odinfit}[2]{%
  \sbox{\odinfigbox}{\ignorespaces#2\unskip}%
  \ifdim\wd\odinfigbox>#1\relax
    \typeout{ODINFIT: figure scaled from \the\wd\odinfigbox\space to \the#1}%
    \resizebox{#1}{!}{\usebox{\odinfigbox}}%
  \else
    \usebox{\odinfigbox}%
  \fi}

% ---- TikZ styles --------------------------------------------------------------------------
\tikzset{
  % every figure
  odin/.style={font=\footnotesize, color=o3ink, line join=round,
               >={Stealth[length=3.4pt, width=2.8pt]},
               % node text breaks the same way in any class: no hyphens, loose lines rather than overfull
               execute at begin node={\hyphenpenalty=10000\exhyphenpenalty=10000%
                 \tolerance=9999\emergencystretch=10pt}},
  % shapes
  box/.style={draw, rectangle, rounded corners=1.6pt, align=center, line width=0.55pt,
              inner xsep=2.5pt, inner ysep=2pt, minimum height=6.5mm},
  store/.style={draw, cylinder, shape border rotate=90, aspect=0.16, align=center,
                line width=0.55pt, inner xsep=2pt, inner ysep=1pt, minimum height=8mm},
  plain/.style={align=center, inner sep=1pt},
  % build status (combine with a shape: [box, built])
  built/.style={fill=o3grn4, draw=o3grn8, solid},
  inprog/.style={fill=o3grn2, draw=o3grn7, solid},
  spec/.style={fill=white, draw=o3grn7, dash pattern=on 2.4pt off 1.4pt},
  planned/.style={fill=white, draw=o3brn5, dash pattern=on 0.6pt off 1.3pt,
                  line width=0.75pt},
  % neutral fills for structure figures
  neutral/.style={fill=o3brn2, draw=o3brn6, solid},
  light/.style={fill=o3grn1, draw=o3grn7, solid},
  external/.style={fill=o3brn2, draw=o3brn6, dash pattern=on 3pt off 1.2pt on 0.6pt off 1.2pt},
  emph/.style={line width=1.3pt},
  % grouped regions, drawn on the background layer
  region/.style={draw=o3brn4, fill=o3brn1, rounded corners=3pt, line width=0.5pt,
                 inner sep=4pt},
  bluregion/.style={draw=o3grn3, fill=o3grn1, rounded corners=3pt, line width=0.5pt,
                    inner sep=4pt},
  regionlabel/.style={font=\scriptsize\bfseries, color=o3brn7, inner sep=1.5pt},
  % arrows
  flow/.style={->, line width=0.55pt, draw=o3ink},
  flowplanned/.style={->, line width=0.6pt, draw=o3brn6, dash pattern=on 0.6pt off 1.3pt},
  ref/.style={->, line width=0.5pt, draw=o3grn8, dash pattern=on 2.4pt off 1.4pt},
  edgelabel/.style={font=\scriptsize, color=o3sub, inner sep=1pt, fill=white},
  % records: a header row and field rows, built with \matrix (see fig_ir_model.tex)
  record/.style={matrix of nodes, ampersand replacement=\&, row sep=-\pgflinewidth,
                 inner sep=0pt, nodes={draw=o3brn6, line width=0.5pt, fill=white,
                 font=\scriptsize, anchor=center, align=left, inner xsep=2.5pt,
                 inner ysep=1.2pt, minimum height=4.4mm}},
  rechead/.style={fill=o3grn3, font=\scriptsize\bfseries, align=center},
  % legend keys
  legendkey/.style={box, minimum width=4.2mm, minimum height=2.6mm, inner sep=0pt},
  legendtext/.style={font=\scriptsize, inner sep=1pt, text height=1.5ex, text depth=0.3ex},
}

% \odinlegend{<node options, e.g. anchor=north east, at={(x,y)}>}
% The build-status key, one row, as a single node so figures can place it like any node.
\newcommand{\odinlegend}[1]{%
  \matrix[ampersand replacement=\&, column sep=1.2pt, inner sep=2pt, draw=o3brn4,
          rounded corners=2pt, fill=white, line width=0.4pt, #1]{%
    \node[legendkey, built]{};   \& \node[legendtext]{built};       \& \node{\hspace{3pt}}; \&
    \node[legendkey, inprog]{};  \& \node[legendtext]{in progress}; \& \node{\hspace{3pt}}; \&
    \node[legendkey, spec]{};    \& \node[legendtext]{specified};   \& \node{\hspace{3pt}}; \&
    \node[legendkey, planned]{}; \& \node[legendtext]{planned};     \\
  };}

% \odinrecord{<name>}{<matrix options: position, nodes={text width=...}>}{<header>}{<rows>}
% A record (struct) as a one-column matrix: a header cell, then one cell per field, rows
% separated by \\ and ended by \\. Cells are <name>-<row>-1, so edges can leave from a field.
\newcommand{\odinrecord}[4]{%
  \matrix[record, #2] (#1) {|[rechead]| #3 \\ #4};}
